\documentclass[12pt]{article}
\usepackage[utf8]{inputenc}
\usepackage[T1]{fontenc}
\usepackage[margin=1in]{geometry}
\usepackage{setspace}
\usepackage{xcolor}
\usepackage[colorlinks=true, urlcolor=blue!60!black, linkcolor=blue!60!black]{hyperref}
\setlength{\parindent}{0em}
\setlength{\parskip}{0.45em}
\renewcommand{\baselinestretch}{1.0}
\begin{document}
\begin{center}
{\Large \textit{Grutter v. Bollinger}, 539 U.S. 306 (2003)\par}
\vspace{0.3em}
{\normalsize Case Brief --- One-Page Edition\par}
\end{center}

\textbf{Memory Jogger.} Diversity is a compelling interest; Michigan Law School's holistic race-conscious admissions upheld 5--4, until \textit{SFFA} (2023) (\href{https://www.courtlistener.com/opinion/130155/grutter-v-bollinger/}{CourtListener}).

\textbf{Facts.} Michigan Law School's 1992 policy, modeled on Powell's \textit{Bakke} opinion, reviewed each file holistically and sought a ``critical mass'' of underrepresented minority students. Barbara Grutter, a white applicant with a 3.8 GPA and a 161 LSAT, was rejected and sued under the Equal Protection Clause, Title VI, and \S~1981. An expert testified that race-blind admissions would cut minority enrollment from 14.5\% to 4\%.

\textbf{Procedural History.} After a bench trial, the district court enjoined the policy; the en banc Sixth Circuit reversed 5--4.

\textbf{Votes.} Affirmed 5--4. O'Connor wrote for the Court; Ginsburg concurred; Scalia and Thomas concurred in part and dissented in part; Rehnquist and Kennedy dissented.

\textbf{Holding.} Diversity is a compelling interest in higher education. Narrow tailoring requires race as a flexible ``plus'' in individualized review, serious consideration of race-neutral alternatives, no undue harm to nonminority applicants, and a limited duration.

\textbf{Reasoning.} The Court deferred to the school's educational judgment and found diversity's benefits ``not theoretical but real''; unlike \textit{Gratz}'s point system, critical mass was no quota, and programs need a ``logical end point.'' Ginsburg would ``hope, but not firmly forecast'' an end in 25 years. Scalia mocked the ``mystical `critical mass' justification.'' Thomas would allow only ``pressing public necessity.'' Rehnquist found admissions tracked applicant shares, proving racial balancing, and Kennedy wrote that ``Deference is antithetical to strict scrutiny, not consistent with it.''

\textbf{Examples.} \emph{Same-Side:} holistic medical school review treating race as one plus factor. \emph{Opposite-Side:} a 15-point racial bonus. \emph{Fence-Sitter:} a ZIP-code preference adopted to preserve diversity.

\textbf{Critique.} Conservatives said deference diluted strict scrutiny, and Richard Sander's contested ``mismatch'' thesis claimed preferences hurt beneficiaries; Derrick Bell called diversity a distraction from remedying subordination.

\textbf{Key Quotation.} ``We expect that 25 years from now, the use of racial preferences will no longer be necessary to further the interest approved today.'' --- Justice O'Connor, 539 U.S. at 343.

\end{document}
